\section{Theoretical Background}
\label{sec:theory}

\subsection{The Jevons Paradox}

In 1865, the British economist William Stanley Jevons observed
that improvements in the efficiency of steam engines did not
reduce coal consumption but instead increased it. As Jevons
argued, ``it is wholly a confusion of ideas to suppose that
the economical use of fuel is equivalent to a diminished
consumption. The very contrary is the truth'' \cite{sorrell2009}.
The underlying mechanism is straightforward: efficiency
improvements lower the effective cost of energy services,
stimulating greater demand and ultimately outweighing the
initial savings. This phenomenon has since been termed the
Jevons Paradox.

\subsection{The Rebound Effect}

The Jevons Paradox forms the theoretical foundation of what is
today referred to as the rebound effect -- an umbrella term for
mechanisms that reduce the potential energy savings from improved
energy efficiency \cite{sorrell2009}. A widely cited illustration
is that of a driver who replaces a vehicle with a more
fuel-efficient model, only to drive further as a result of the
lower running costs \cite{sorrell2009}. Research suggests that
rebound effects are consistently larger than conventional
assessments assume \cite{sorrell2009}.

\subsection{Applying the Rebound Effect to Artificial Intelligence}

The rebound effect is not limited to conventional energy systems.
As AI technologies make logistics, industrial production, and
data-intensive computing more efficient, the resulting cost
reductions tend to expand usage rather than reduce overall
resource consumption \cite{mhlanga2025}. The fact that an AI
model becomes more efficient per query does not imply that
overall resource consumption will decrease \cite{luccioni2025}.
AI also generates indirect impacts through behavioral and
structural changes that affect other processes and lifestyles
\cite{luccioni2025}. The following sections examine both
dimensions in the specific context of AI deployment.
